% Generated draft: requires mathematical and source review.
\documentclass[11pt]{article}
\usepackage[margin=25mm]{geometry}
\usepackage{fontspec}
\setmainfont{Latin Modern Roman}
\usepackage{amsmath,amssymb,mathtools,unicode-math}
\setmathfont{Latin Modern Math}
\usepackage{enumitem,needspace,xurl,hyperref}
\setlength{\parindent}{0pt}
\setlength{\parskip}{5pt}
\setlength{\emergencystretch}{3em}
\newcommand{\N}{\mathbb{N}}
\newcommand{\Z}{\mathbb{Z}}
\newcommand{\Q}{\mathbb{Q}}
\newcommand{\R}{\mathbb{R}}
\newcommand{\C}{\mathbb{C}}
\newcommand{\F}{\mathbb{F}}
\newcommand{\A}{\mathbb{A}}
\newcommand{\PP}{\mathbb{P}}
\newcommand{\E}{\mathbb{E}}
\newcommand{\eps}{\varepsilon}
\newcommand{\dd}{\,\mathrm{d}}
\begin{document}
\section*{UnsolvedMath ID: 30003137}
\textbf{Source problem code:} catalogue entry
\subsection{Definitions and mathematical statement}
For integers \(n\geq 1\), let \([n]=\{1,\ldots,n\}\). For an integer \(k\geq 3\), a nontrivial \(k\)-term arithmetic progression in \([n]\) is a set
\[
\{a,a+d,\ldots,a+(k-1)d\},
\qquad a,d\in\Z,\quad a\geq1,\quad d\geq1,\quad a+(k-1)d\leq n.
\]
Here the conventional interpretation of arithmetic progressions with distinct terms is used; the supplied statement does not explicitly specify nontriviality. A subset of \([n]\) is \(k\)-term-AP-free if it contains no such progression. Define
\[
\mathcal F_k(n)=\{A\subseteq[n]:A\text{ is }k\text{-term-AP-free}\},
\qquad
r_k(n)=\max_{A\in\mathcal F_k(n)}|A|.
\]
The enumeration includes every such subset, including the empty set.

The question is whether, for every fixed integer \(k\geq3\),
\[
|\mathcal F_k(n)|=2^{(1+o(1))r_k(n)}
\qquad\text{as }n\to\infty
\]
holds along positive integers \(n\). Equivalently, is
\[
\lim_{n\to\infty}\frac{\log_2|\mathcal F_k(n)|}{r_k(n)}=1
\qquad\text{for every fixed integer }k\geq3?
\]
The error term may depend on \(k\); no uniformity for growing \(k\) is requested.

The supplied record labels this question open and its statement \texttt{corrected\_verified}; these are supplied curation labels, not independent verification. Its literature assessment concerns bounds for \(r_3(n)\), rather than directly establishing the status of this enumeration question. The report-year discrepancy described in the Sources also remains unresolved.
\subsection{Short English statement}
For each fixed integer \(k\geq3\), count all subsets of \(\{1,\ldots,n\}\) containing no nontrivial \(k\)-term arithmetic progression. Is the base-two logarithm of this count asymptotic to the largest possible size \(r_k(n)\) of such a subset? The supplied open status requires independent review.
\subsection{Sources}
\begin{itemize}
\item Local UnsolvedMath registry, supplied record with permanent ID \(30003137\), source problem code \texttt{catalogue entry}, title \emph{Asymptotic Enumeration of Arithmetic-Progression-Free Sets}, in the Oberwolfach Reports Open Problems collection. The record gives proposed year 2016 and names no individual proposer.
\item Supplied report citation: \emph{Combinatorics and Probability (2017)}, Oberwolfach Reports, DOI 10.4171/owr/2016/22, pp.\ 1189--1257. \url{https://doi.org/10.4171/owr/2016/22}. The supplied background describes Oberwolfach Reports as documenting mathematical workshops and their open-problem discussions. The citation's 2017 label differs from the 2016 proposed year and the year embedded in the DOI; this discrepancy has not been resolved.
\item Supplied statement-verification note: the statement was checked against the official Oberwolfach report and preserves its hypotheses, notation, and open mathematical target. This is a claim in the supplied record; no independent source consultation is asserted here.
\item Kelley and Meka, \emph{Strong Bounds for 3-Progressions}. \url{https://arxiv.org/abs/2302.05537}.
\item Bloom and Sisask, \emph{Improvement to the Kelley--Meka Bounds}. \url{https://arxiv.org/abs/2309.02353}.
\item Literature assessment dated exactly \texttt{2026-08-21}: supplied dated curation, not independently verified. It labels the status open and states that the exact growth of \(r_3(n)\) remains unknown despite major improvements, with Kelley--Meka and Bloom--Sisask upper bounds still not matching Behrend-type lower bounds. The background explicitly describes this as dated literature triage derived from the supplied OWR curation and requests independent verification of time-sensitive status and cited claims. This assessment addresses the extremal function and does not itself verify the status of the enumeration question.
\end{itemize}
\end{document}
